% Copyright (c) 2026 Zhongpai Gao. Released under Apache-2.0.
% Native analytic source and continuation theory: Hailey Collet.
% Two-level returned-source construction: ArviFoxy. Inherited notices retained.
\documentclass[10pt,a4paper]{article}
\usepackage[margin=19mm]{geometry}
\usepackage{amsmath,amssymb,amsthm,mathrsfs,booktabs,microtype}
\usepackage[colorlinks=true,urlcolor=blue,citecolor=blue,linkcolor=blue]{hyperref}
\setlength{\parindent}{0pt}\setlength{\parskip}{3.5pt}
\setlength{\emergencystretch}{2em}
\newcommand{\N}{\mathrm N}\newcommand{\E}{\mathcal E}
\newcommand{\tht}{\vartheta}\newcommand{\kap}{\kappa}
\newcommand{\cS}{\mathcal S}\newcommand{\cD}{\mathcal D}
\newcommand{\pos}[1]{(#1)_+}
\newtheorem{theorem}{Theorem}\newtheorem{lemma}{Lemma}
\newtheorem{proposition}{Proposition}
\title{Finite-depth affine source returns for Dirichlet \texorpdfstring{$L$}{L}-functions}
\author{Zhongpai Gao}\date{20 October 2026}
\begin{document}\maketitle\vspace{-2em}
\begin{abstract}
We describe a five-stage affine source construction that retains the actual
residual reduction through recursive LOW steps and counts the returned HIGH
records jointly over their complete arithmetic histories. The source weight
keeps every intermediate row scale; common reduction is credited only at its
native valuation-cone rate. Exact rational certificates give a positive
budget for all later HIGH bins. Backward induction closes the uniform
whole-root moment, and the inherited detector and weighted-numerator method
gives the all-moduli Dirichlet nonvanishing boundary $0.8748087$.
The accompanying Lean proof constructs the moment from the actual finite
sources and checks the resulting unconditional theorem.
\end{abstract}

\section{The initial moment and the nonvanishing theorem}
Work in the Eisenstein Hecke framework of OpenAI and Collet
\cite{openai,collet}. Its global zero supremum is denoted by $\beta$.
The native input price is $\kap_0=4499/6000$; it is an established input,
independent of the conclusion. Put
\[
 \tht=\frac{8748087}{10^7},\qquad
 \lambda=3.741331371862391\ldots,\quad
 B=0.982029506308942\ldots,\quad s_+=0.04038098.       \tag{1}
\]
The exact rational $\lambda,B$ are specified in the accompanying parameter
file. All displayed decimal junctions and widths below are exact terminating
rationals, except the two explicitly truncated values in (1).

The source construction gives the following moment. Let $Z>1$ be the root
radius, $m_i=\log_Z X_i$, and $s=\sum_j w_j$ the selected prime width.
Write $\E_\eta(\Phi,\mathbf p)$ for the literal initial energy: its original
prime pools, annular coefficients, fixed mask, fixed $Q$, and seed one are
those of the inherited initial-moment theorem. The radial control is
$\cD(\Phi)$ and the finite profile control is $\cS_{\mathscr S}(\mathbf p)$.

\begin{proposition}[Whole-root moment]\label{prop:moment}
Suppose $51/100\le\beta\le7/8$. For every $\varepsilon>0$ a positive
prime-slot mesh is chosen before the ambient cardinality and weighted source.
For each fixed original weighted source and admissible $\eta$ unchanged by
excluding the unit ideal, fix its inherited $Q$ and modulus/support data.
There are a finite seminorm set
$\mathscr S$, a degree $D$, and $C>0$, chosen before $Z$ and the moving
profiles, such that eventually
\[
 \E_\eta(\Phi,\mathbf p)
 \le C\cD(\Phi)\cS_{\mathscr S}(\mathbf p)^2
       (1+|t|+T)^D Z^{1+\varepsilon},                 \tag{M}
\]
for all shifts bounded by $T$, both reflection orientations, and all actual
initial sources satisfying
\[
 \frac{81}{200}\le m_i\le\frac{264}{625},\qquad
 \frac{17}{500}\le s\le s_+,\qquad m_1+m_2+\lambda s\le B.
\]
The profile supports lie in $[1/4,9/4]$ and the original annular exponent is
$33/50$. The same proved estimate supplies the detector and conjugated
numerator, with their original source data.\end{proposition}

\begin{theorem}[Nonvanishing]\label{thm:endpoint}
For every $n\ge1$, every
Dirichlet character $\chi$ modulo $n$, and every $z\in\mathbb C$,
\[
 \Re z>\tht,\qquad \neg(\chi=1\text{ and }z=1)
 \quad\Longrightarrow\quad L(\chi,z)\ne0.
\]
\end{theorem}
ArviFoxy's two-level construction \cite{arvi} returns its second-step HIGH
part instead of estimating it independently. We extend that organization to
one root step and four inner stages. The new issue is preserving a single
actual source, its residual factor, and its raw weight across every depth.
The following two sections give the corresponding payment and count.

\newpage
\section{The actual source, its residual factor, and its raw scale}
Normalize logarithms by $\log Z$. At one physical source let $C_1,C_2$ be
the two actual common ideals, with frozen nonzero allocation coefficients.
For the selected residual ideal $L$ and allocation $a$, retain
\[
 F=\log_Z\!\bigl(\operatorname{formalReductionFactor}(L,a)\bigr)\ge0.
\]
If $A$ is the original log-volume, $z$ the selected width, and $u$ the
retained nominal width, the native supported-source identities give
\begin{align}
 n_1+n_2+z'&\le A-\log_Z\N C_1-\log_Z\N C_2-F+P,\notag\\
 n_1+n_2+6\kap_0z'-u
 &\le G+6a_0+\omega+P-H_C-F.                         \tag{2}
\end{align}
Here $P$ is the measured support padding, $G$ the parent's old-capacity
defect, $a_0$ its amplifier, and $H_C$ the actual common-norm payment after
the nominal-width subtraction. Both inequalities concern the \emph{same}
$L,a$ and the same supported source; (2) follows by subtracting the two raw
allocation reductions from the original volume and then taking logarithms.

This preservation is compatible with source selection. For a finite positive
source sum, $E\le\sum_i c_iE_i$ with $c_i\ge0$, choose an index maximizing
$E_i/B_i$, where $B_i>0$ is its original reciprocal/Cauchy normalization.
Then
\[
 E\le\Bigl(\sum_i c_iB_i\Bigr)\frac{E_{i_*}}{B_{i_*}}. \tag{3}
\]
Thus (2) and (3) retain the same source and residual factor. Finite physical
choices retain actual provenance. When heights vary, a factor-two supremum
selects one whole licensed arithmetic majorant with its full path and count
receipts; it need not be an attained analytic run. Terminal depths are not
maximized separately.

\begin{lemma}[Raw attachment]\label{lem:raw}
At stage $j$ let $U_j$ be the physical unit, $Z_j=Z^{U_j}$ the local base,
$m_j$ the signed first-mass dyad loss, and $X_0,\ldots,X_3$ the native
coordinates in the original base. Set $\widehat X_0=X_0+U_jm_j$.
If the LOW edge deletes actual width $D_j$, its raw exponent satisfies
\[
 e_j=M_j+\widehat X_0+q_jD_j+e_{j+1}+\epsilon_j,\qquad
 M_{j+1}+\log_Z\N(v_{j+1})-\log_Z\N(v_j)\le X_2+P_j. \tag{4}
\]
Here $v_j$ is the actual prefix multiplier, $M_j$ the actual row scale, and
$\epsilon_j$ contains the remaining paid finite fees. The historical missing
radical $X_1$ enters the count with coefficient one.
\end{lemma}
\begin{proof}
The source contains $\widehat X_0+a_j+q_jD_j$; the AtState factor contains
$M_j-a_j$. Thus $a_j$ cancels and $M_j$ remains. The mass term is not
assumed small. The \emph{same} first bin supplies a side lower bound $b_j$
and label-side exponent $S_j$ with
\[
 m_j=3\delta_j-b_j,\qquad m_j+S_j\le3\delta_j+6\log_{Z_j}2.
\]
This is the native shifted-coordinate payment: the large mass loss cancels
its matching label side before iteration. The actual multiplier identity
gives the relative-scale inequality in (4). Substitute it and
$L_{j+1}\le X_3-D_j-e_{\mathrm{vol}}+P'_j$ in the child payment to obtain
\[
 (e_j-M_j)+(X_1+\mathrm{miss}_{j+1})
       +\rho(q)\bigl(r_{j+1}-\log_Z\N v_j\bigr).
\]
There is no reciprocal cancellation of the unnormalized recursive LOW term.
The volume credit is $q e_{\mathrm{vol}}$; common reduction already in $X_3$
keeps only its native valuation-cone credit.
\end{proof}

A volume majorant requires compensation. For native local volume $A$, a lift
by $\Delta$ changes $(A,m,X_0,X_3)$ to
$(A+\Delta,m-\Delta,X_0+U\Delta,X_3+U\Delta)$.
Hence $\widehat X_0$, and therefore the complete raw edge weight, is unchanged;
the label side increases by $\Delta$ and its cancellation also survives.
The lifted $m$ is a compensated coordinate, not a new mass estimate.
At the root $U_0=1$ and $\Delta=A_0-\widetilde A_0\ge0$, where
$\widetilde A_0$ is the actual volume and $A_0$ its root-cell majorant.
The volume coordinates below use this compensated lift, not an unpriced
replacement of the actual volume.

A narrow, exhausted, old-capacity or fully reflected child returns an empty
LOW-continuation forest using the native direct/old bound. A continuing child
uses the next actual region. One genuine small/coarse height splice preserves
the return linearly. Root radius $Z$, seed one and row scale one follow from
the initial source reduction; the fixed conductor's logarithmic fit loss is
absorbed after its eventual threshold.

\newpage
\section{Five-stage payment and a single terminal count}
For $q\le3/4$, put
\[
 \rho(q)=1-\frac{10q}{3(5-2q)},\quad c(q)=5-6\rho(q),\quad
 h(q)=\max\{3/19,\,4\rho(q)-5/4\}.
\]
The five junctions are
\[
 (q_0,q_1,q_2,q_3,q_4)=
 (.638021915,.676506091,.693781289,.711716695,.730301837). \tag{5}
\]
Starting with $(k_4,I_4)=(\kap_0,1)$, the preceding output region is defined
recursively, for $j=3,2,1,0$, by
\[
 K=\frac{6(k_{j+1}-q_{j+1}I_{j+1})}{q_{j+1}c(q_{j+1})},\quad
 6k_j-1=\frac{6k_{j+1}-1}{1+Kq_{j+1}},\quad
 I_j=\frac{I_{j+1}+K(1-\rho(q_{j+1}))-10^{-8}}{1+Kq_{j+1}}. \tag{6}
\]
The inner affine slopes are $6k_{j-1}$. Their first-error channels use prime
fraction $99/100$; the terminal step uses the original price $\kap_0$.

For an actual path, let $U_j,z_j$ be its retained and slot widths, and
$A_j$ its compensated volume coordinate, all in the original base. Write $V_j\ge0$ for missing-width loss,
$R_j\ge0$ for the augmented common/residual reduction and $\ell_j\ge0$
for additional slot loss. The source gives, up to explicitly paid rounding,
\[
 U_{j+1}=U_j-6a_j-V_j,\quad A_{j+1}=A_j-d_j-V_j-R_j,\quad
 z_{j+1}=z_j-d_j-\ell_j,\qquad \ell_j\le V_j+R_j.       \tag{7}
\]
Here $d_j$ is the model deletion and $a_j$ the affine amplifier; the actual
detector deletion and finite mesh surcharge remain in (4).
Let $V=\sum_{i<j}V_i$, $R=\sum_{i<j}R_i$,
$a=\sum_{i\le j}a_i$, $P=\sum_{i<j}q_id_i$,
$W=1-V-6a$ and $L=A_0-V-R-\sum_{i<j}d_i$, where $A_0$ is the
uniform root-cell upper bound for the original volume.
The terminal payment licensed by the actual source is
\[
 J_{\rm paid}(q)=(1-\rho(q))W-qL+a-P
          +\bigl(1-\rho(q)-h(q)\bigr)V-h(q)R.           \tag{8}
\]
The $-h(q)R$ correction prevents counting common reduction twice.
For $q\ge q_j$, the uncorrected first four terms have increment
\[
 \left(\frac{(50/3)W}{(5-2q)(5-2q_j)}-L\right)(q-q_j).\tag{9}
\]
Here $W\ge0$. Both correction coefficients in (8) are nondecreasing in $q$;
a junction certificate and $\frac{50W}{3(5-2q_j)^2}-L\ge0$ pay every later bin.
The exact certificates give $J_{\rm paid}(q)>10^{-7}$ for each returned
inner depth, including the paid coordinate errors; the root own-HIGH part
uses its separate marked estimate. Each affine certificate is an
identity $f=\gamma+\sum_i t_i g_i$, $t_i\ge0$: when $g_i\ge-E$ it gives
$f\ge\gamma-E\sum_i t_i$. This also pays the finite moving grid.

The terminal HIGH estimate is applied once, over complete histories. A key
retains both common-label pairs, their inactive subsets, the native seed,
the four dyads and the orientation. If
$m(H)=\sum_{i<j}X_{1,i}$ is the historical missing exponent, the prefix-fibre
count is $\ll_\epsilon Z^{m+\epsilon}$ at a fixed multiplier ideal.
Squarefree-divisor bounds pay the seed and fixed dyad boxes pay the remaining
multiplicity. The residual $L,a,F$ remains source metadata, not a free count
coordinate. Thus the native zero-density estimate gives
\[
 \sum_H w_H\,\mathrm{HIGH}_H
 \ll C K(1+T)^D\!\prod_{i<j}(|\mathcal B_i|+1)Z^{E+\epsilon}, \tag{10}
\]
provided $w_H Z^{a_\sigma+b_\sigma+\rho(q_\sigma)r_\sigma+m(H)}\le Z^E$
for every native class $\sigma$. Here $a_\sigma,b_\sigma,r_\sigma$ are its
original weighted, missing and row exponents and $q_\sigma=2\beta_\sigma-1$.
Equations (4), (8) supply this payment with the original $w_H$.
The zero-height exponent can be chosen before the box budget, class cap and
fixed analytic return coefficient; these affect only the eventual threshold.

\newpage
\section{Closing the moment and applying the detector}
The source construction is finite. At the last inner stage the native LOW
estimate is terminal. At each preceding stage the actual LOW source produces
one whole licensed returned forest, while the local HIGH source produces its
own native record. The old allowance is summed with its original first-seed
reciprocal; the returned forest is not divided by that seed norm. Residual
projector shells, both ray sums and positive buffered bins keep the original
normalization and introduce only their fixed constants and subpower losses.
The three-height splice is paid once. Finite unions of seminorms and maxima
of degrees precede every moving source.

\begin{proof}[Proof of Proposition~\ref{prop:moment}]
The terminal source starts backward induction through the four inner stages.
Each preceding step invokes the actual canonical, OLD, coarse and physical
sources on the same node; the resulting complete receipt supplies the next
step. All primitive constants and thresholds are chosen over the finite
original subsets and moving ledger grid before $Z$, uniformly under rebasing.
The root LOW source then invokes this proved induction and the joint count;
its marked HIGH estimate supplies the separate root record. The initial
source reduction and a final coarse-height extension give (M). Thus the
literal root-moment interface is constructed from the source proof.
\end{proof}

\begin{proof}[Proof of Theorem~\ref{thm:endpoint}]
Suppose $\beta>\tht$. The inherited unconditional bound $\beta\le7/8$
and $\tht>51/100$ admit Proposition~\ref{prop:moment}. Apply it to the
actual detector dictionary and conjugated weighted-numerator source. Both
consumers prove the required excluded-unit character identity. Finite
dictionary unions give common seminorms, degrees and thresholds before $Z$.
Keep the outer detector density $\delta$ distinct from the inner HIGH $q$.

For inverse length $r$, plain length $m$, and amplitude $p$, the inverse,
affine-plain and long-witness counts have exponents
\[
 1-\delta\bigl(p+(1-p)r\bigr),\quad
 1-2\delta m-2\delta p\frac{B-2m}{\lambda},\quad
 1-\delta+(5/6-\delta)(r+m-1).                          \tag{11}
\]
Use the affine count only when its actual widths lie in (M)'s window; the
capped-width and plain-top branches cover the adjacent ranges, and the
inherited count covers the complement. The inverse exponent $5/6$ is
unchanged. The affine balance, for example, is
\[
 R_{\rm aff}=1-\delta+
 \frac{(5/6-\delta)\delta(1-p)(1-2pB/\lambda)}
 {(5/6-\delta)(3-p-4p/\lambda)+\delta(1-p)(2-4p/\lambda)}.
\]
Set $h=0.79716116522$, $\ell=11/3-4\tht$, $y=h-2\ell$ and
$\kap_*=2\tht-1$. The old and affine capacities are
\[
 C_O(d)=\frac{\min(d,2y-d)}{6\kap_*},\qquad
 C_N(d)=\frac{2y-(2-B)d}{\lambda}.
\]
The inherited weighted comparison, with its selected count exponent $R$, is
\[
 \frac{1+\delta}{2}+\frac{13h}{75}-\frac y2-\frac\ell2+\ell\delta p
 +d\left(\frac14-\frac{17}{50}\right)-\frac{\delta pC}{2}
 +\frac{3dR}{4}\ \le\ \tht-\frac1{50000000}.           \tag{12}
\]
The rational endpoint certificates cover $0\le\delta\le3/4$,
$0\le p\le1/2$, with positive denominators and all source admissions.
They use the old-capacity ends and the affine-capacity junction and top;
affinity in $d$ supplies the intervals. Clearing denominators leaves
polynomials expressed in the nonnegative tensor Bernstein basis, so their
positive coefficients prove (12) on each closed rectangle.

The standard feedback comparison extends (12) from
$\kap_*$ to the actual $2\beta-1$. Choose the remaining positive losses
inside its strict reserve before the target characters and heights.
Levy's weighted numerator and Collet's continuation theorem
\cite{levy,collet} then improve the same global zero supremum, a contradiction.
Hence $\beta\le\tht$. The inherited Hecke-to-Dirichlet transfer gives
the conclusion, with the principal pole excluded.
\end{proof}

\small\textbf{Acknowledgments and notices.}
The analytic framework is due to OpenAI; the affine source and continuation
theory to Hailey Collet; the weighted numerator to Akash Levy; asymmetric
marked counting to timbalud; the affine detector to DaysSky; and the two-level
returned-source construction to ArviFoxy. Earlier balanced marked ledgers
and related affine counting are also used. Zhongpai Gao used AI assistance.
Copyright and license notices of all inherited sources are retained.
This manuscript is released under Apache-2.0, subject to those notices.

\begin{thebibliography}{9}\footnotesize
\bibitem{openai} OpenAI, \emph{Quasi-Riemann hypothesis formalization},
\href{https://github.com/openai/math/blob/fd4aeeb2ee4fc729c18d98444fed42fd0529eeeb/lean/docs/003.md}{framework and manuscript}, snapshot \texttt{fd4aeeb2}.
\bibitem{collet} Hailey Collet, affine moment and continuation theory,
\href{https://github.com/GTimG/QuasiRiemannTracker/pull/42}{PR42}.
\bibitem{levy} Akash Levy, weighted numerator,
\href{https://github.com/GTimG/QuasiRiemannTracker/pull/6}{PR6}.
\bibitem{arvi} ArviFoxy, \emph{A second affine step for the LOW children},
\href{https://github.com/GTimG/QuasiRiemannTracker/pull/67}{PR67}.
\bibitem{marked} timbalud, asymmetric marked counts,
\href{https://github.com/GTimG/QuasiRiemannTracker/pull/50}{PR50}; DaysSky, affine detector,
\href{https://github.com/GTimG/QuasiRiemannTracker/pull/56}{PR56}.
\bibitem{balanced} Zhongpai Gao, balanced common geometry and adaptive marked source,
\href{https://github.com/GTimG/QuasiRiemannTracker/pull/63}{PR63}.
\end{thebibliography}
\end{document}
